\subsection{WPIC background controls and time-history diagnostics}
\label{sec:wpic-diagnostics}

The weighted Vlasov--Poisson and subhalo tasks share one background and one
marker evolution. The options \code{--vp-universe}, \code{--vp-alpha}, and
\code{--vp-time-offset} select that background before the initial conditions,
field, or particle trajectories are constructed. Let \(t\) be elapsed run
time, \(t_0>0\) the time offset, \(A(t)\) the scale factor, \(H(t)\) its
fractional expansion rate, and \(\alpha\) the prescribed exponent. For the
expanding choice,
\[
 A(t)=\left(\frac{t+t_0}{t_0}\right)^\alpha,\qquad
 H(t)=\frac{\alpha}{t+t_0},\qquad
 \bar\rho(t)=\frac{3H(t)^2}{8\pi G}.                    \tag{WPIC-D1}
\]
The defaults are \(\alpha=2/3\) and \(t_0=12\). The driver admits finite
\(0\le\alpha\le1\); the prescribed background and the source convention
must match the physical problem. The static choice sets \(\alpha=0\),
\(A=1\), \(H=0\), and \(\bar\rho=0\). Its supplied cusp is an initial
density distribution whose equilibrium properties require independent checks.
The mesh, collision law, and adaptive collision clocks retain their existing
controls.

\codesource{Shared WPIC background selection}{\path{cpu/src/drivers/expanding_wpic.cpp}}{Construct one consistent background for the initial conditions, field, and canonical-comoving particle push}
\begin{lstlisting}[language=C++]
if (result.vp_universe == "static") result.vp_alpha = 0.0;
// make_parameters uses these resolved values:
p.vp_universe = run.vp_universe;
p.vp_alpha = run.vp_alpha;
p.vp_time_offset = run.vp_time_offset;
\end{lstlisting}

A fresh run writes \path{background_contract.txt}. Restarting requires the
same background options as that file. Earlier checkpoints with the prescribed
expanding \(\alpha=2/3,t_0=12\) background can be imported with those values.
The separate \path{radial_mesh_contract.txt} checks the mesh. Keeping both
records beside the marker checkpoint preserves the interpretation of saved
positions and velocities.

The default static sample gives most particles tangential velocities close
to their circular-orbit values. Reducing these velocities initiates a
controlled collapse. The option \code{--initial-velocity-factor} selects a positive
factor \(F\), with default \(F=1\). For a fresh sample it changes physical
velocity about the halo center according to
\[
 u_{r,\mathrm{new}}=F(Hr+w),\qquad
 w_{\mathrm{new}}=F(Hr+w)-Hr,\qquad
 \eta_{\mathrm{new}}=F^2\eta.                         \tag{WPIC-D1a}
\]
Positions and represented masses retain their sampled values. In a static
background the initial kinetic energy changes by \(F^2\), while the
initial gravitational field retains its value. Thus \(F=1/4\) prepares a
halo with one sixteenth of the default kinetic energy. In an expanding
background the same prescription scales physical velocity, with the saved
peculiar velocity obtained by subtracting the original Hubble flow.

\codesource{Controlled initial velocities}{\path{cpu/src/drivers/expanding_wpic.cpp}}{Scale fresh physical velocities once while preserving the density sample}
\begin{lstlisting}[language=C++]
marker.w = run.initial_velocity_factor
    * (marker.w + hubble_velocity) - hubble_velocity;
marker.eta *= run.initial_velocity_factor
    * run.initial_velocity_factor;
\end{lstlisting}

The additional \path{initial_condition_contract.txt} records \(\gamma_i\)
and \(F\). A continuation repeats those options and evolves its saved
velocities directly. The factor acts once during fresh initialization.
Legacy checkpoints without this contract retain the factor-one convention.

Sachiet Kapur's \path{cpu/tools/plot_wpic_history.py} reads
\path{history.csv} and makes five aligned panels. They show the maximum
resolved cell density, the mobile and field matter masses, the stored energy,
the applied and limiting timesteps, and the minimum effective marker count.
The density maximum may occur away from the center. The mobile-mass column
includes deposited mass and the unresolved inner carrier. The field-mass
column uses the Poisson solver's matter-mass integral. Compare these masses with the declared Poisson source and boundary
treatment. Matter inside the unresolved inner region belongs to the mobile
population while its contribution to the resolved profile depends on the
field convention.

The stored energy diagnostic is
\[
 E_{\rm hist}=\sum_{p\ {\rm active}}\left[
 \frac{m_p}{2}\left(w_p^2+\frac{2\eta_p}{r_p^2}
                              +\Phi(r_p)\right)+U_p\right]. \tag{WPIC-D2}
\]
Here \(m_p\) is represented marker mass, \(w_p\) is peculiar radial velocity,
\(\eta_p=L_p^2/2\), \(\Phi\) is the self-consistent potential, and \(U_p\)
is the internal energy carried by a subhalo marker. Each gravitational pair contributes to the potential at both partners.
The factor one half on \(\Phi\) therefore counts its energy once. Expansion, the unresolved inner source,
and boundary work enter the complete physical energy budget. The annotation
reports the change of the stored series; numerical conservation is assessed
with the budget appropriate to that run.

\codesource{WPIC history report}{\path{cpu/tools/plot_wpic_history.py}}{Density, mass, energy, timestep, and sampling histories from saved output}
\begin{lstlisting}[language=bash]
python3 cpu/tools/plot_wpic_history.py --history RUN/history.csv \
  --output output/wpic_history.png
\end{lstlisting}

\subsubsection{Density profiles and inner and outer cusp indices}

The tool
\path{cpu/tools/render_weighted_pic_gamma_fits_gifs.py} generates
initial/final density and slope comparisons,
inner/outer index histories, a CSV of the fitted values, fit settings in JSON,
and an exact frame-time table. It reads \code{gamma_i} and marker sampling
from the run metadata. Complete snapshots are checked before frame selection,
and the first and last complete times are retained at every stride.

For positive radius \(r\) and density \(\rho\), define
\[
 x=\log_{10}r,\qquad y=\log_{10}\rho,\qquad
 s=\frac{\dd y}{\dd x},\qquad \gamma=-s.               \tag{WPIC-D3}
\]
The derivative panel displays the signed slope \(s\). The time-history
panels display \(\gamma_{\rm in}\) and \(\gamma_{\rm out}\), so a declining
profile \(\rho\propto r^{-3/2}\) has \(s=-3/2\) and \(\gamma=3/2\).
The derivative uses a local quadratic fit with a 21-point window; estimates
within half a window of either edge are omitted. The local quadratic uses neighboring samples to distinguish the slope from
curvature. A meaningful estimate therefore requires enough resolved points
across the fitting window.

The inner straight-line fit uses a default center \(x=-0.75\) and half-width
\(0.25\). The two command-line fit controls select these values explicitly.
The outer fit searches inward from the outermost cell with effective marker
count at least 200, locating the first point where the absolute second
logarithmic derivative is below 2 or above 6. Its center is the midpoint in
\(x\) between that feature and the sampled outer edge; its half-width is
\(0.25\). Its fit uses cells satisfying the marker-count cut. These numerical
choices define an empirical profile-feature estimate. A physical outer-cusp
interpretation also needs inspection of the fit interval, comparison of
neighboring windows, and marker and mesh convergence. An unavailable fit is
saved as an undefined value in the CSV. The density floor is a display cut;
the simulated density and mass remain the input data.

\codesource{Extended density and cusp-index report}{\path{cpu/tools/render_weighted_pic_gamma_fits_gifs.py}}{Read saved profiles and retain differentiation, fit, and frame-selection settings}
\begin{lstlisting}[language=bash]
python3 cpu/tools/render_weighted_pic_gamma_fits_gifs.py \
  --run-dir RUN --output-directory output/density \
  --stride 10 --inner-fit-center-log10r -0.75 \
  --inner-fit-half-width-log10r 0.25 --shoulder-neff-min 200
\end{lstlisting}

\subsubsection{Resolved Newtonian virial ratio}

Sachiet's \path{cpu/tools/plot_subhalo_virial_1.py} forms a Newtonian report
from complete logarithmic-mesh profiles. VENUS stores arithmetic radial
centers \(r_i=(r_{i-1/2}+r_{i+1/2})/2\). For consecutive face ratio
\(q=r_{i+1}/r_i\), the corresponding faces and shell masses are
\[
 r_{i-1/2}=\frac{2r_i}{1+q},\qquad
 r_{i+1/2}=q r_{i-1/2},\qquad
 \Delta m_i=\frac{4\pi\rho_i}{3}
                  (r_{i+1/2}^3-r_{i-1/2}^3).           \tag{WPIC-D4}
\]
This reconstructs the simulation volumes, including cells with zero density.
Let \(\bar w_i\) be the mean radial peculiar velocity,
\(\sigma_{r,i}^2\) its variance, and \(\langle v_t^2\rangle_i\) the sum of
the two tangential velocity second moments. With \(G=1\), the report uses
\[
 T=\frac12\sum_i\Delta m_i
       (\sigma_{r,i}^2+\bar w_i^2+\langle v_t^2\rangle_i),\qquad
 W=-\sum_i\frac{\Delta m_i}{r_i}
       \left(\sum_{j<i}\Delta m_j+\frac{\Delta m_i}{2}\right),\qquad
 Q=\frac{2T}{|W|}.                                     \tag{WPIC-D5}
\]
The report also writes a CSV with time, virial ratio, kinetic and potential
energies, and resolved mass. When the metadata supplies the represented
mobile mass, an additional column gives the resolved fraction of that mass.
This fraction identifies when matter carried inside the unresolved inner
region materially changes the radial domain represented by the ratio.
Complete snapshots are selected before applying the requested stride.

The half-shell mass in \(W\) approximates each shell's contribution to its
own gravitational energy. This diagnostic combines resolved orbital motion
with resolved self-gravity. Expansion,
boundary pressure, the unresolved inner carrier, and subhalo internal energy
have separate contributions in a complete virial budget. A relativistic halo
is assessed with the Einstein--Vlasov balance and time-series criteria in the
EF chapter. The reference \(Q=1\) supplies the isolated Newtonian comparison
under its equilibrium assumptions.

The orbital equations show explicitly how expansion enters this comparison.
For a fixed set of material particles between encounter or escape events,
define its radial second moment \(I=\sum_p m_pr_p^2\) and mixed moment
\(C=\sum_p m_pr_pw_p\). A dot denotes a physical-time derivative.
Substituting \(\dot r=Hr+w\) and
\(\dot w=L^2/r^3-\Phi_r-Hw\) gives
\[
 \dot I=2HI+2C,\qquad
 \dot C=2T+W_\Phi,\qquad
 W_\Phi=-\sum_p m_pr_p\Phi_r(r_p).                    \tag{WPIC-D6}
\]
Here \(T\) includes radial and tangential peculiar motion of that material
set. The centrifugal term in \(\dot C\) supplies its tangential kinetic
energy. Differentiating the first relation and substituting the second gives
\[
 \frac12\ddot I-H\dot I-\dot H I=2T+W_\Phi.           \tag{WPIC-D7}
\]
The force moment \(W_\Phi\) includes the declared background and central
gravity sources. Its definition therefore differs from the resolved-shell
self-gravity quadrature \(W\) in (WPIC-D5). A fixed radial observation volume
also has surface transport and pressure terms, and a capture or a change in
the deposited carrier has its own moment change. For a closed isolated
stationary Newtonian system, \(H=0\), \(\ddot I=0\), and the self-gravity
force moment equals its potential energy. Equation (WPIC-D7) then yields
\(2T+W=0\). Sustained settling of the measured ratio, density profile, and
the applicable moment and boundary budgets supplies the physical evidence
for virialization.

The driver also writes \path{virial_budget.csv} at each history output.
It evaluates \(I,C,T\), and \(W_\Phi\) over every active marker, including
matter carried inside the resolved mesh. The force sample is the same
interpolated field used by the orbit update. Alongside these quantities it
records resolved and hidden mass, physical kinetic energy
\(T_{\rm phys}=T+HC+H^2I/2\), internal energy, \(H\), and \(\dot H\).
Its columns \code{C_dot_orbit} and \code{I_ddot_half_orbit} give the
instantaneous orbital right-hand sides of the moment equations. A finite
difference through a collision or reflecting-wall event also includes the
corresponding jump or impulse.

\codesource{Applied-force virial moments}{\path{cpu/src/drivers/expanding_wpic.cpp}}{Include the same gravity as the marker pusher in the global orbital balance}
\begin{lstlisting}[language=C++]
result.inertia += mass * r * r;
result.radial_moment += mass * r * w;
result.kinetic_peculiar += 0.5L * mass * (w * w + vt2);
result.kinetic_physical += 0.5L * mass * (u * u + vt2);
result.force_virial -= mass * r * sample.phi_r;
\end{lstlisting}

An unresolved central mass \(M_{\rm h}\) attracts the resolved halo even
when its mass lies outside the profile integral. Its approximate cross
contribution to the resolved force moment is
\[
 W_{\rm cross}=-G M_{\rm h}\sum_i\frac{\Delta m_i}{r_i}. \tag{WPIC-D8}
\]
Comparing this term with (WPIC-D5) identifies the effect of the omitted inner
source. The all-marker budget additionally retains inner kinetic energy
and the actual force interpolation. Numerical refinement then tests the
discretization of the same physical balance.

\codesource{Resolved Newtonian virial report}{\path{cpu/tools/plot_subhalo_virial_1.py}}{Integrate saved shell moments with the simulation's arithmetic-center geometry}
\begin{lstlisting}[language=bash]
python3 cpu/tools/plot_subhalo_virial_1.py --run-dir RUN \
  --output output/newtonian_virial.png --stride 10
\end{lstlisting}

\subsubsection{Initial data, time-colored profiles, and radius scaling}

The tool \path{cpu/tools/plot_initial.py} reads the state actually written
by a zero-time grid-solver setup. It follows the model name in the metadata
when selecting relativistic metric panels. A comparison of initial conditions
therefore uses the same density, velocity moments, and phase-space output
that the chosen solver will evolve.

\codesource{Initial-state visualization}{\path{cpu/tools/plot_initial.py}}{Inspect the grid solver's saved initial profiles and phase-space distribution}
\begin{lstlisting}[language=bash]
python3 cpu/tools/plot_initial.py --input-dir SETUP_OUTPUT \
  --output-dir output/initial_state
\end{lstlisting}

The tool \path{cpu/tools/plot_color_rho_r.py} overlays saved WPIC density
profiles, using color to identify physical time. The tool
\path{cpu/tools/plot_wpic_scaled_profiles.py} additionally selects a radius
\(R_f(t)\) enclosing a chosen fraction \(0<f<1\) of the resolved mass.
It integrates the shell masses in (WPIC-D4). Within a shell of constant cell
average density, the enclosed mass is linear in \(r^3\), which determines
\(R_f\) by interpolation. Its radius-scaled panel shows
\(r/R_f\); the optional density normalization shows
\(\rho R_f^3/M_{\rm resolved}\). The complete resolved mesh determines
\(R_f\) before the effective-marker display cut is applied. A self-similar profile has a time-independent shape and amplitude in these
scaled coordinates. Compare both over several times, while accounting
separately for the inner carrier and boundaries.

\codesource{Radius-scaled WPIC profiles}{\path{cpu/tools/plot_wpic_scaled_profiles.py}}{Compare resolved profiles using measured enclosed-mass radii and optional density scaling}
\begin{lstlisting}[language=bash]
python3 cpu/tools/plot_wpic_scaled_profiles.py 'RUN/profile_*.csv' \
  --num 6 --frac 0.1 --normalize-density --output output/scaled_profiles.png
\end{lstlisting}

\subsubsection{Radial dispersion and mean flow}

The material-WPIC profile writer uses explicit columns
\code{radial_mean} and \code{radial_variance}; the latter is already the
centered variance.  Its \code{tangential_second} is the sum of the two
tangential second moments.  Persistent-marker SphDSMC instead writes
\(j=\rho\overline v_r\), centered pressure \(p=\rho\sigma_r^2\), and
\(p_t=\rho\langle v_t^2\rangle\).  These conventions determine which
mean-flow term enters the kinetic energy.

The radial dispersion measures motion about the local mean velocity.
Spherical DSMC writes its pressure column as \(p=\rho\sigma_r^2\),
while the Newtonian grid writes \(p=\rho\langle v_r^2\rangle\).
Both write radial momentum density \(j=\rho\langle v_r\rangle\).
The renderer \path{cpu/tools/render_density_slope_sigma2_gif.py} therefore uses
\[
 \sigma_r^2=\begin{cases}
 p/\rho, &\text{spherical DSMC pressure},\\
 p/\rho-(j/\rho)^2, &\text{Newtonian grid second moment}.
 \end{cases}                                             \tag{WPIC-D9}
\]
The metadata field \code{bp_collision_operator=spherical_dsmc} selects the
first convention. Other Newtonian tables use the second convention;
\code{--radial-moment-convention central} or \code{raw} specifies it for
external tables. The selected input convention is saved beside the GIF.

\codesource{Radial pressure convention}{\path{cpu/tools/render_density_slope_sigma2_gif.py}}{Convert the raw second moment to central pressure before combining radial cells}
\begin{lstlisting}[language=Python]
if convention == "raw":
    profiles["p"] = (
        profiles["rho"]
        * radial_dispersion_squared(profiles, "raw")
    )
\end{lstlisting}
Combining radial cells must preserve their mass, momentum, and kinetic
energy. Adaptive bins therefore first sum mass, radial momentum, and the full
second velocity moment. Subtracting the merged mean velocity squared then retains
both the original dispersions and differences of bulk velocity between
cells. This permits a physical comparison of an outer dispersion bump with
the accompanying mean radial flow.

Explicitly, merging cells of mass \(M_i\), radial mean \(u_i\), and centered
variance \(\sigma_i^2\) gives
\[
 \sigma_{\rm merged}^2=
 \frac{\sum_i M_i(\sigma_i^2+u_i^2)}{\sum_i M_i}
 -\left(\frac{\sum_iM_i u_i}{\sum_iM_i}\right)^2.        \tag{WPIC-D10}
\]
The first term retains the kinetic energy carried by differences in mean
flow between the cells; the second removes the merged bin's own mean flow.

The Frank benchmark renderers additionally apply fixed unit conversions.
Their stated convention is \(G=M_{\rm total}=r_s=1\), with
\(M_{\rm total}/M_s=1.1252580933331495\).  A native
\(G=M_s=r_s=1\) profile uses a different density, velocity, and time
conversion.  Establish that convention from the initial-data record before
using a benchmark renderer for a physical comparison.

\subsubsection{Native initial states and reflecting-wall measurements}

A native particle file can supply a different initial density distribution.
The CPU driver accepts \code{--initial-condition-label gaussian_shell} with
an explicit \code{--restart} input. This label describes the supplied
particles. The checkpoint checksum identifies their actual state. For this
input the metadata gives the actual total marker mass and the mass inside
the first resolved face, together with \code{native_input_sha256}.
The physical radial velocity is \(u=Hr+w\), with \(H=0\) in a static run.
Retaining the same \(r,u,\eta,m\) and particle identities therefore requires
\(w=u-Hr\) in the expanding initial file; here \(m\) is represented mass
and \(\eta=L^2/2\) for specific angular momentum \(L\).

\codesource{Describe a native Gaussian initial state}
{\path{cpu/src/drivers/expanding_wpic.cpp}}
{Evolve the supplied particles and record their actual mass and checksum}
\begin{lstlisting}[language=bash]
cpu/venus_cpu_vlasov_poisson_wpic_expanding \
  --restart INITIAL/restart_state.bin \
  --initial-condition-label gaussian_shell --vp-universe static \
  --gamma 0.5 --initial-velocity-factor 1 \
  --radial-cells 512 --radial-minimum 0.01 --radial-maximum 10000 \
  --maximum-timestep 0.015 --end-time 10000 \
  --boundary-diagnostics --output-directory RUN
\end{lstlisting}
The background, mesh and initial-condition sidecars remain beside the native
input. The descriptive label is independent of those existing restart
contracts. A native state at \(t=0\) also writes its initial profile when
this custom label is supplied. The profile can then be checked against the
intended Gaussian density before inspecting its evolution.

The optional \code{--boundary-diagnostics} counts specular impacts at the
outer wall. It also records the largest particle radius among accepted
step endpoints. If \(n_k\) is the number of wall impacts in accepted step
\(k\), the segment budget and endpoint radius are
\[
 N_{\rm wall}=\sum_{k\,\rm accepted}n_k,\qquad
 r_{\rm max,end}=\max_{k,p}r_p(t_k).                 \tag{WPIC-D11}
\]
Here \(p\) labels particles and \(t_k\) is an accepted endpoint. Regular
passages through the origin preserve the particle trajectory and contribute
zero wall impacts. A trial step rejected by the collision-clock check is
rolled back before either diagnostic enters the budget.

\codesource{Accepted-step wall budget}
{\path{cpu/src/drivers/expanding_wpic.cpp}}
{Accumulate physical wall impacts after the orbit step is accepted}
\begin{lstlisting}[language=C++]
segment_wall_impacts +=
    orbit_step_diagnostics.reflecting_wall_impacts;
\end{lstlisting}
The file \path{boundary_history.csv} gives the accepted step and time,
segment start step and time, cumulative outer impacts, maximum endpoint
radius and its recorded endpoint time, and the physical outer-wall radius.
Profiles, checkpoints and scheduler stops flush this record. Each process
continuation starts an explicitly identified segment; its final cumulative
count supplies that segment's contribution. These measurements establish
whether the chosen physical wall is reached during the calculation.

\subsubsection{Complete particles and their self-consistent potential}

The GPU production driver can retain selected particle states together with
the field reconstructed from those same particles. The schedule is empty by
default. Every selected time lies on the existing profile schedule, so the
output uses the established orbit landing times. The callback follows the
completed KDK kick and any due collision batch.

\codesource{Select complete particle and field output}
{\path{gpu/src/wpic_production_driver.cu}}
{Add selected snapshots to an otherwise unchanged GPU invocation}
\begin{lstlisting}[language=bash]
--output-interval 10 \
--phase-space-snapshot-times 0,10,30,100,300,800,2000,5000,10000
\end{lstlisting}
Each \path{phase_space/snapshot_NNNNN} directory contains the native
\path{particles.bin}, actual \path{field.json}, and joint
\path{manifest.json}. The particle file retains \(m,r,w,\eta\), identities,
internal states and restart clocks in the portable CPU/GPU format. The field
file retains the potential \(\Phi\), its radial derivative, enclosed source
masses, line density, mesh, moments, inner mass and expansion parameters.
It states whether particle weights describe total density or density
contrast, and records the solver's zero-at-infinity potential convention.
The manifest gives actual and requested time, accepted step, both file
checksums, input checkpoint and executable checksums, and the full invocation.
The complete group is installed with one directory rename.

\codesource{Serialize an existing endpoint field}
{\path{gpu/src/wpic_production_driver.cu}}
{Save the field already reconstructed at the profile callback}
\begin{lstlisting}[language=C++]
venus_gpu::driver::write_phase_space_field_json(
    (temporary / "field.json").string(), restart.time,
    field, mesh, options,
    run.universe == UniverseMode::eds_expanding);
\end{lstlisting}
An exact continuation preserves an existing snapshot whose particle and
field checksums agree. Existing snapshots later than the supplied restart
cause a chronology error. A continuation in a fresh output directory saves
the selected states reached from its input time. The CPU production driver
continues to write its normal native checkpoints; the selected paired output
above belongs to the GPU production driver. The portable particle file can
be read by either native checkpoint reader.

To reconstruct the spherical distribution, integrate over the position and
tangential-velocity angles. Their factors \(4\pi\) and \(2\pi\), together
with \(v_t\,dv_t=d\eta/r^2\), give
\[
 d\Gamma=8\pi^2\,dr\,du\,d\eta,\qquad
 f_{\rm bin}=\frac{\sum_{p\in\rm bin}m_p}
 {8\pi^2\,\Delta r\,\Delta u\,\Delta\eta}.           \tag{WPIC-D12}
\]
Here \(d\Gamma\) is physical phase volume, \(v_t\) is tangential speed,
\(f_{\rm bin}\) is the bin-averaged mass distribution function, and each
\(\Delta\) denotes that bin's width. The change from \(w\) to \(u=w+Hr\)
has unit Jacobian at fixed radius and time. A projection over \(\eta\)
gives mass per \(dr\,du\). For a logarithmic radius axis, retain the actual
\(\Delta r\), or explicitly normalize by \(\Delta\log_{10}r\).

The same saved field gives each particle's instantaneous orbital energy per
unit mass,
\[
 \epsilon_p=\frac{u_p^2}{2}+\frac{\eta_p}{r_p^2}
                  +\Phi(r_p,t).                    \tag{WPIC-D13}
\]
Use the solver's constant-line-mass reconstruction for \(\Phi(r_p,t)\).
The centrifugal term has its regular zero limit for \(L=0\) at the origin.
A weighted histogram then gives \(dM/d\epsilon\), retaining the full mass
at both negative and positive orbital energies. In an expanding run
\(\Phi\) is the peculiar potential and \(\epsilon_p\) is an instantaneous
halo binding diagnostic. Subsequent motion and the boundary budget determine
escape. An energy-dependent phase-space density additionally requires the
phase volume of each energy bin. The reconstruction and these conventions
are described in \path{gpu/docs/PHASE_SPACE_SNAPSHOTS.md}.
